\chapter*{Conclusion}
\label{chap:conclusion}
\addcontentsline{toc}{chapter}{\textbf{Conclusion}}

Ce mémoire a pour objectif de concevoir un outil d'aide à la décision permettant de détecter la fraude en assurance automobile et d'orienter les efforts d'investigation vers les dossiers les plus à risque. Il vise ainsi à répondre à une difficulté structurelle de la lutte anti-fraude : le coût d'une investigation manuelle, qui rend non rentable le contrôle de la fraude de faible montant et laisse échapper une part importante du préjudice. L'enjeu était donc de montrer qu'un agent d'intelligence artificielle pouvait automatiser cette détection à moindre coût et à grande échelle.

Pour y parvenir, il a d'abord fallu constituer une base de travail réaliste. Faute de données réelles suffisamment riches, nous avons enrichi une base publique de sinistres automobiles de vingt-neuf variables simulées, reproduisant l'information dont dispose réellement un assureur au moment d'un sinistre (historique de l'assuré, du réparateur, du véhicule et du tiers, comportement déclaratif, etc.). Ce choix, assumé comme une limite, a permis de rapprocher l'exercice des conditions opérationnelles réelles.

La démarche s'est ensuite structurée autour de la construction d'un agent piloté par un \textit{prompt} décomposé en étapes successives, reproduisant la logique actuarielle : ingestion et audit de la qualité des données, analyse statistique descriptive, puis construction et comparaison des modèles de prédiction de fraude. Trois modèles ont été calibrés puis évalués selon un arbitrage entre interprétabilité et performance : un \textit{scoring} en points, un arbre de décision et un modèle de Gradient \textit{Boosting}.

Ce dernier s'est révélé le plus discriminant, sans toutefois disqualifier le \textit{scoring} en points, qui demeure une alternative crédible lorsque la transparence prime. Nous avons par ailleurs constaté que les trois modèles faisaient converger leur pouvoir explicatif vers les mêmes variables clés, ce qui a renforcé la confiance dans la cohérence des résultats avec l'intuition métier.

\textit{FAIRE UN PARAGRAPHE SUR LE SECOND OUTIL PARCE QU'ON EN PARLE SANS L'AVOIR INTRODUIT}

L'un des partis pris forts de ce travail a été de décliner l'outil en deux modules indépendants, pensés pour deux profils d'utilisateurs distincts. Le premier, destiné aux actuaires, prend en charge la construction et la calibration du modèle ; le second, destiné aux gestionnaires de sinistres, en permet l'exploitation quotidienne sans en exposer la complexité, en restituant pour chaque dossier une probabilité de fraude et une orientation d'investigation. Cette séparation des rôles répond à une réalité de terrain : ceux qui conçoivent le modèle et ceux qui l'utilisent ne font pas parti des mêmes équipes.

Enfin, l'outil a été pensé non comme un modèle figé mais comme un dispositif évolutif : une boucle de rétroaction réintègre dans la base d'apprentissage les dossiers dont l'investigation est close, permettant un recalibrage régulier sur des cas avérés, tandis qu'un rapport de pilotage mesure en continu la valeur ajoutée de l'outil (fraude évitée, impact sur le ratio Sinistres/Primes, retour sur investissement).

Ce travail comporte naturellement des limites. La principale tient au caractère partiellement simulé des données et à la sur-représentation de la fraude dans la base retenue, qui imposeraient de recalibrer le modèle sur une sinistralité réellement représentative avant tout déploiement. Par ailleurs, l'approche supervisée adoptée ne détecte que la fraude de type connu, celle dont on possède des exemples historiques labellisés ; elle reste par construction aveugle aux schémas de fraude organisée ou inédits. Plusieurs prolongements permettraient de dépasser ces limites : le couplage à une base portefeuille, l'intégration de données externes, le recours à des méthodes non supervisées ou d'analyse de graphes pour la fraude en réseau, la mutualisation inter-assureurs pour repérer les fraudeurs « nomades », ou encore l'extension de l'outil à d'autres branches de l'assurance non-vie.

Au-delà de la seule performance technique, ce mémoire illustre la place croissante de l'actuaire dans la conception d'outils décisionnels concrets, à l'intersection de la modélisation, de la contrainte réglementaire et de l'usage métier. La détection de la fraude est devenue un enjeu stratégique pour les assureurs, et si l'intelligence artificielle en démultiplie aujourd'hui les moyens, elle ne prendra tout son sens qu'articulée à une supervision humaine garante de la maîtrise du risque et du respect des droits des assurés.